\documentclass[script,con]{debate}
\motion{「勇敢的人先享受世界」是一个陷阱 / 不是一个陷阱}
\stance{反方}{不是一个陷阱}
\meta{赛制}{中国科学技术大学新生辩论赛（反方后立论，反四先结辩）}
\meta{口径来源}{备赛文档 第十节 10.2 反方口径表。定义、判准、论点、数据均取自该表；允许换说法，不允许换意思、换数字。}
\begin{document}
\debatetitle
\begin{thesisbox}[全场一句话立论]它给的是一个顺序，不是一张支票；说好的「更早开始」兑现了，没兑现的那部分是准备的事，不是这九个字的事。\end{thesisbox}
\begin{thesisbox}[全场只守两句话]① 我方承认陷阱可以没有设局者，我方要的是那层伪装。② 全场只问一句：这句话盖住的那个东西，到底是哪几个字？\end{thesisbox}

\begin{speech}{1. 反一开篇立论稿}{正文 755 字 / 180 秒}
谢谢主席。开场我先把一件事让给对方：我方承认，陷阱可以没有设局者。中等收入陷阱没有人挖，它照样叫陷阱。今天我方不会问对方「是谁设的局」，那样问是在躲比赛。

我方接受对方的判准，也接受那三件事：一，外观和真实的因果对不对得上；二，照做的人受损是不是常态；三，这份损失是不是被它遮住的东西造成的。请评委注意一个不对称：三件缺一，它就不是陷阱。对方要立三件，我方只要打掉一件。今天我方打第一件和第三件。

我方第一点。它给的是一个顺序，不是一张支票。陷阱要有一层伪装，伪装要有一个被盖住的落差，也就是说好的和实际给的之间那个差。那么这句话说好的是什么？它说的是，在先动和先等之间，先动的人更早开始。请评委想一想，这件事有没有兑现？兑现了。它从来没说你会赚到什么，会留住什么，会不会中途回头。二〇二三年这句话最早只是国庆假期旅游九宫格的一句配文，没有签名，没有对价，更没有一个收款的人在后面等着。智联招聘二〇二五年春季三万七千多人的调研里，三成五有跳槽意愿的人会选择裸辞，其中两成一说的是「休息后再出发」。这是一群心里有数的人，不是一群被九个字骗进坑的人。说白了，他们知道自己在换什么，只是对方不愿意承认这一点。

我方第二点，账要记在准备上。让人摔的是没留退路，不是有勇气。勇敢这个词，词典里说的是有勇气有胆量，而勇气针对的是\stress{已知的}风险；不知道有风险还往前冲，汉语里叫莽撞。Bandura 一九七七年讲过，一个人最靠得住的能力感，来自亲手做成过的事。所以行动本身产出的是真东西。至于后悔，一九九四年 Gilovich 的研究说人更后悔没做的事，可二〇二三年近两千六百人的公开复制里，长期是五成一对四成九。我方承认这一点，而且我方要的正是这个五五开：做和不做，长期看后悔的机会一样大。一个五五开的选择，可以叫冒险，可以叫赌，唯独不叫陷阱。陷阱是进去就亏。

我方全场只问对方一句话：这句话盖住的那个东西，到底是哪几个字？答不出来，伪装这一格就是空的。

我方从不说这句话有多好。说它不严谨可以，说它片面也可以。但不严谨的话和圈套，隔着整整三件事。谢谢。
\end{speech}
\begin{contingency}
\ifthen{正四质询已拿到我方某个承认}{开头第一句先钉住：「刚才我方承认的是陷阱不需要设局者，不是承认这句话骗了人。」}
\ifthen{对方把我方打成「替这句话辩护」}{立刻回：「我方从没说它好。我方说的是它不满足陷阱的要件。」}
\ifthen{时间不够}{删掉配文来源那一句，保留智联的两个数字。}
\end{contingency}

\begin{stage}{2. 反一被质询防守预案}{正四质询，120 秒双边计时}
\textbf{原则}：双边计时，每答不超过八秒。承认该承认的，绝不说「必须有设局者」。被逼二选一时先指出预设，再给第三种说法。
\begin{qa}
\qarow{中等收入陷阱有设局者吗？}{没有。我方承认陷阱不需要设局者。}
\qarow{那你方还剩什么理由？}{伪装。陷阱要有一个被盖住的落差，请对方交出来。}
\qarow{这句话承诺过你会过得更好吗？}{没有。它承诺的是先动的人更早开始。}
\qarow{那它凭什么让人辞职？}{它提供的是一个排序，不是一个保证。人照着排序做决定，还要自己算账。}
\qarow{「享受」两个字你方怎么解释？}{享受是结果的名字，不是结果的保证。原句没有说享受得到多少、留多久。}
\qarow{这九个字里哪个字提到了存款？}{一个都没有。格言本来就不列条件，这不叫伪装。}
\qarow{那你承认它省略了条件？}{承认。省略是不严谨，不是圈套。}
\qarow{勇敢和鲁莽你方怎么分？}{勇气针对已知的风险，不知道有风险还冲的叫莽撞。}
\qarow{是不是只要准备好就不叫勇敢了？}{不是。准备好之后仍然可能失败，那时候上才叫勇敢。}
\qarow{你方是不是在替消费主义说话？}{不是。我方一句都没夸过这句话，我方只说它不是陷阱。}
\end{qa}
\end{stage}

\begin{stage}{3. 反四质询正一问题链}{120 秒，双边计时}
\textbf{总目标}：逼正方交出「被盖住的那个东西」，并拿到「正方没有受损数据」这个承认。
\begin{chain}{链一：伪装链}{逼正方指出被盖住的内容}
\q{陷阱是不是要有一层伪装？}{答「是」：进 2；答「不是」：问「那和错话有何区别」}
\q{伪装是不是要盖住点什么？}{答「是」：进 3}
\q{这句话盖住的是哪几个字？}{答不出即收束}
\end{chain}
\begin{chain}{链二：受损链}{拿到「没有直接数据」的承认}
\q{你方说受损是常态，有数据吗？}{答「没有」：进 2；答「有」：问「哪一份」}
\q{那常态两个字，凭什么成立？}{答「靠机制」：进 3}
\q{机制能代替常态的举证吗？}{不等答完，进链三}
\end{chain}
\begin{chain}{链三：边界链}{逼正方划出「什么不是陷阱」}
\q{「努力就有回报」是不是陷阱？}{答「是」：记下来；答「不是」：进 2}
\q{它和这句话的区别在哪？}{答「顶替」：进 3}
\q{那顶替有没有证据，还是你方的读法？}{收}
\end{chain}
\cut{好的，我理解为您方没有正面回答。下一个问题。}
\close{感谢。对方到现在没有说出被盖住的是哪几个字，也承认了没有受损数据，我方二辩会继续。}
\end{stage}

\begin{speech}{4. 反二驳论稿}{正文 505 字 / 120 秒}
对方立论听上去很有力，可它有一个洞，而且是致命的：\stress{对方全场只证明了这句话不准确，没有证明它是个坑。}

第一，判准这一格对方欠着账。对方自己立的三件事，第二件是「照做的人受损是常态」。可对方一辩到现在，没有拿出任何一条关于「照它做的人后来怎么样了」的证据。对方唯一的数据，是四成多人不敢 gap。请评委看清楚，这条数据说的是这些人\stress{没有走}。没走的人怎么会掉进坑里？对方拿一份「人们清醒地算过账」的数据，去论证「人们被外观骗了」，这不是举证，这是反证。

第二，对方说这句话把门槛换成了勇气。可这九个字从头到尾没有提过任何条件。你想想，哪一句格言会把前提条件列在后面？我方承认，它不严谨。但请对方回答一个问题：省略和顶替，对方怎么区分？对方说「早睡早起身体好」没有顶替，这句话有。凭什么？凭对方自己的读法。原句只说了勇敢的人怎么样，它没有说不勇敢的人是谁。那句「你不够勇敢」，是对方替它补上去的。

第三，对方引一九四三年的幸存者偏差。这一点我方完全同意，可幸存者偏差是所有成功叙事的共同属性。学长的经验帖、课本上的科学家故事、今天晚上任何一场分享会，全都只有活着回来的那一架飞机。按对方的标准，它们全是陷阱。一个把人类全部经验传递都判成陷阱的标准，衡量的已经不是陷阱了，它衡量的是这个世界有没有例外。

回到我方判准：三件缺一，它就不是陷阱。对方今天缺的是第二件。谢谢。
\end{speech}
\begin{contingency}
\ifthen{正一没有给出「顶替」这条线}{把第二驳点扩成 60 秒，直接追问「省略和顶替怎么分」。}
\ifthen{正一主打商业收编}{加一句：「那对方打的是广告主，不是这九个字，请把被告席上的人换对。」}
\end{contingency}

\begin{stage}{5. 反二对辩要点}{90 秒，单边计时}
\begin{points}{进攻点（4–5 个，每个 15–20 秒，说完就停）}
\item 被盖住的是哪几个字？请给四个字。
\item 您方那四成多人是没走的人，没走的人怎么掉进坑？
\item 省略和顶替，您方拿什么分？除了您方的读法还有别的吗？
\item 幸存者偏差人人都有，按您方标准学长的经验帖也是陷阱。
\item 受损是常态，这是您方立的第二件事，数据在哪？
\end{points}
\begin{points}{备用点（6–8 个，用于回应）}
\item 我方承认它不严谨，不严谨和圈套隔着三件事。
\item 我方也承认陷阱不需要设局者，所以请不要再打定义杀这一格。
\item 五五开是您方给的数字，五五开的选择不叫陷阱。
\item 原句没有说不勇敢的人是谁，那句判决是您方补的。
\item 勇气针对已知风险，词典写的，不是我方加的。
\item 不含在校生的青年失业率一成五左右，还在连续回落。
\item 智联那份报告测的是意愿，您方也不能拿它说结果。
\item 您方越强调要算账，越是在说准备，那正是我方第二点。
\end{points}
\begin{points}{对方最可能问的三个问题 → 我方回应}
\item 「享受两个字你方删哪去了？」→ 没删。享受是结果的名字，不是结果的保证，原句没说享受多少、留多久。
\item 「一句什么都没说的话，怎么会红三年？」→ 我方从没说它什么都没说。它说的是别把一生花在等条件齐了。
\item 「你方是不是在替消费主义说话？」→ 我方一句都没夸过它。我方说的是它不满足陷阱的要件。
\end{points}
\end{stage}

\begin{stage}{6. 反三盘问问题链}{90 秒，单边计时}
\textbf{赛前确认}：向组委会确认被盘问方回答是否计时；若不计时，每条链可多加一问。
\begin{chain}{链一：常态链（先问正二，再问正四）}{检验两人对「受损是常态」的答案是否一致}
\q{（问正二）受损是常态，这一件立住了吗？}{答「立住了」：进 2；答「靠机制」：进 3}
\q{（问正四）四辩，同样的问题，您方有数据吗？}{两人不一致就当场点名，小结重锤}
\q{（问正二）那第二件靠什么成立，机制还是数据？}{答「机制」：指出举证责任没交}
\end{chain}
\begin{chain}{链二：边界链}{逼正方划线，看这条线能不能守住}
\q{（问正四）「努力就有回报」是不是陷阱？}{答「不是」：进 2；答「是」：记下来}
\q{（问正四）区别是顶替，那顶替有证据吗？}{答「有」：问「哪一条」；答「是读法」：收}
\q{（问正二）二辩，您也认为那是读法吗？}{两人答案不同即收束}
\end{chain}
\textbf{盘问目标}：小结里要证明两件事——正方的第二件事至今没有证据；正方所谓「顶替」，只是正方自己的一种读法。
\end{stage}

\begin{stage}{7. 反二、反四被盘问防守预案}{两人口径必须一致}
\begin{qaa}
\qaarow{陷阱需不需要有设局者？}{不需要。我方一辩已经承认了。}{不需要。这一格我方不打。}
\qaarow{那你方还剩什么？}{伪装和归因两格。}{伪装和归因，请对方交出被盖住的东西。}
\qaarow{「享受」两个字怎么解释？}{享受是结果的名字，不是结果的保证。}{原句没说享受得到多少、留多久。}
\qaarow{你方承不承认它省略了条件？}{承认。格言本来就不列条件。}{承认。省略是不严谨，不是圈套。}
\qaarow{你方承不承认有人因此摔了？}{承认有人摔了。摔的是没留退路的人。}{承认个案存在。个案不是常态，常态要对方举证。}
\qaarow{你方是不是在说这句话很好？}{我方一句都没夸过它。}{我方从不评价它好不好，只判断它是不是陷阱。}
\end{qaa}
\end{stage}

\begin{speech}{8. 反三盘问小结稿}{正文 317 字 + 临场位 75 字 / 90 秒}
刚才的盘问里，对方交出了两样东西。

第一样，对方承认，「受损是常态」这一件事，靠的是机制，不是数据。可这一件是对方自己立的，也是对方自己说要全扛的。举证责任没交，这一格就是空的。我方承认，这句话确实不严谨；可不严谨的话满大街都是，它们不都是坑。

第二样，对方承认所谓「顶替」，是对方的一种读法。各位评委，这很关键。原句只说了勇敢的人怎么样，它没有一个字说不勇敢的人是谁。那句「你走不了是因为你不够勇敢」，是对方替它补上去的，然后对方再拿这句补出来的话，来证明它是陷阱。

所以到现在，双方真正的分歧是这个：一句省略了条件的话，和一个盖住了条件的坑，差别在哪。我方的答案是差在有没有那层伪装，而伪装总得盖住点什么。一句把话说短了的格言，和一个把东西藏起来的坑，不是一回事。请问，这个东西在哪里？对方被追了四次，到现在一个字也没有交出来。

\reserve{约 80 字，回应正三小结。若正三重锤「改写」，回：改写的是对方，原句里「不勇敢的人」五个字从来不存在。若正三重锤「设局者」，回：那一格我方一辩就让了，请对方不要再打一个我方没有的靶子。}
\end{speech}

\begin{stage}{9. 自由辩战场设计}{240 秒}
\begin{battleground}{战场一：伪装（前 90 秒，主战场）}
\begin{fields}
\field{目标}{让评委看到，正方始终交不出「被盖住的那个东西」。}
\field{开场问题}{这句话盖住的是哪几个字？}
\field{对方答「兜底能力」→ 追问}{哪句话盖住了它？格言不列条件，这叫省略。请问省略和盖住怎么分？}
\field{对方答「它把条件改名叫勇敢」→ 追问}{改名这个动作，是原句做的，还是您方读出来的？请指出是哪个字。}
\field{对方可能反攻 → 我方防守}{对方会说「享受两个字就是承诺」。回：享受是结果的名字，不是结果的保证。原句没说享受得到多少、留多久。}
\field{收束语}{四次追问，对方给不出被盖住的那几个字。伪装这一格是空的，空的一格撑不起一个陷阱。}
\end{fields}
\end{battleground}
\begin{battleground}{战场二：常态（中 90 秒）}
\begin{fields}
\field{目标}{让评委记住，正方自己立的第二件事，全场没有证据。}
\field{开场问题}{照它做的人普遍变差，数据在哪？}
\field{对方答「没有直接数据，这句话才流行三年」→ 追问}{那「常态」两个字从哪来？没有证据的常态，是形容词。}
\field{对方答「常态就是代价错配」→ 追问}{代价错配是您方的描述，不是观察。请问观察在哪？}
\field{对方可能反攻 → 我方防守}{对方会说我方也没有反面数据。回：举证责任在您方，这是您方自己在一辩稿里说全扛的。}
\field{收束语}{对方立了三件事，第二件到现在是空的。三件缺一，它就只是一句不够严谨的话。}
\end{fields}
\end{battleground}
\begin{points}{必问（前两分钟内问完）}
\item 这句话盖住的是哪几个字？
\item 「受损是常态」，数据在哪？
\item 省略和顶替，除了您方的读法还有什么分得开？
\item 「努力就有回报」是不是陷阱？请给一个字的回答。
\end{points}
\begin{points}{必答（15 秒以内正面回答）}
\item 「享受两个字你方删哪去了？」→ 没删。享受是结果的名字，不是保证；原句没说享受得到多少、留多久。
\item 「一句什么都没说的话怎么会红三年？」→ 我方从没说它什么都没说，它说的是别把一生花在等条件齐了。
\item 「你方承不承认有人因此摔了？」→ 承认。摔的是没留退路的人，那是准备的事。
\end{points}
\textbf{时间分配与站位}：前 90 秒战场一，四个必问全部抛完；中 90 秒战场二；最后 60 秒只重复「被盖住的是哪几个字」，不开新战场，留 20 秒备用。一辩守定义与判准，二辩接驳论过的点，三辩接盘问成果，四辩收束与价值。
\end{stage}

\begin{speech}{10. 反四结辩稿}{正文 885 字 / 210 秒，封路式}
谢谢主席。我先做一件事：把正方四辩等一下要走的路，提前铺给各位看。

正方四辩等一下多半会讲一个故事。故事里有一个晚上，有一个人，他看着手机里别人的极光，然后对自己说「我不够勇敢」。这个故事会很动人，我也相信真的有这样的人。但请评委在听的时候，同时问自己一个问题：那句「我不够勇敢」，是这九个字说的，还是他自己说的？原句只说了勇敢的人怎么样，它从头到尾没有一个字说不勇敢的人是谁。正方全场做的，就是把这句没有说出口的话补上去，然后拿它当罪证。

正方四辩还会说，一个需要被改写才守得住的立论，本身就说明原句有问题。这句话很漂亮，可请各位回想全场：被改写的到底是谁的版本？是我方把「享受」读成了「开始」，还是正方把原句读成了「你不配，因为你不够勇敢」？

现在我把账结一下。

双方共同接受了三件事：有没有伪装，受损是不是常态，账该记给谁。正方在一辩稿里明确说过，这三件他们全扛。

第一件，伪装。伪装要有一个被盖住的东西。我方问了四次，从质询问到自由辩：这句话盖住的是哪几个字？正方给不出。格言本来就不列条件，我方承认它不严谨，可不严谨叫省略，不叫盖住。省略是把话说短了，盖住是把东西藏起来，这中间差着一个动作。

第二件，受损是常态。正方亲口承认，没有直接数据。正方唯一的数据是四成多人不敢 gap，可那是没有走的人。一份证明人们清醒地算过账的数据，正方拿来证明人们被外观骗了。

第三件，归因。正方说不出这九个字里哪一个字让人别算账，因为一个都没有。

三件事，正方立住了一件半。而按正方自己在一辩稿里说的判准，缺一件，它就不成立。

最后，我想说说我为什么在意这道题。

我认识一个人，他在一个岗位上待了六年。他不是不敢辞职，他是每年都跟自己说，再等一等，等存款再多一点，等年底奖金发了，等这个项目结束。去年他四十岁。他跟我说的原话是，我这辈子好像一直在等一个从来不会到的时刻。

今天这道题，我方从头到尾没有夸过那九个字。它片面，它省略，它把复杂的事说得太轻。这些我方全认。

可当我们把一句劝人别一直等下去的话，定性成一个陷阱的时候，我们其实是在给那个等了六年的人，递上第七个不动的理由。而且这一次，这个理由听上去比以前所有理由都体面，因为它叫清醒。

清醒是好东西。我方要的也是清醒。但清醒该用来看清代价，不该用来给自己上一把更结实的锁。

一句话可以不严谨，可以片面，可以被人用坏。但它不是一个坑。我方认为，「勇敢的人先享受世界」不是一个陷阱。谢谢。
\end{speech}
\begin{contingency}
\ifthen{正三小结已经重锤「改写」}{开头封路段直接改成回应：「正方三辩说我方改写了原句，可原句里『不勇敢的人』这五个字，从来不存在。」}
\ifthen{正方全场没有给出受损数据}{把第二件事的篇幅加倍，其余各减一句。}
\ifthen{正方在自由辩交出了「被盖住的是兜底能力」}{封路段加一句：「请评委记住，正方直到自由辩才第一次回答这个问题，而且答案是格言本来就不会写的东西。」}
\end{contingency}
\end{document}
